\section{Vertex contributions}

There are also vertex diagrams
involving gluons that connect the gauge link with the gluon lines, see \mbox{Fig. \ref{link}. }

We use the method of calculation
described in Ref. \cite{Balitsky:1987bk}. It is based on the background-field technique,
with the gluon propagator taken
in the ``background-Feynman'' (bF) gauge \cite{Balitsky:1987bk}.
It should be noted that the three-gluon vertex in the bF gauge is different
from the usual Yang-Mills vertex (see e.g. \cite{Abbott:1980hw}).
Therefore, the results obtained for separate diagrams in the
bF gauge differ from those obtained in the usual Feynman gauge
and only the sum of all diagrams must be the same.

\subsection{UV divergent term}

Clearly, the gluon exchange produces a correction just to one of the fields in the
$G_{\mu\alpha}(z)G_{\lambda \beta }(0)$ operator, while another remains intact.
In particular, the diagram \ref{link}a changes $G_{\mu\alpha}(z)$ into the sum of two
terms. One of them contains UV divergences, while the other one is UV finite.
\begin{figure}[t]
\centerline{\includegraphics[width=2.5in]{images/vertgluonno}}
\vspace{-5mm}
\caption{Vertex diagrams with gluons coming out of the gauge link.
\label{link}}
\end{figure}

The UV-divergent term is given by
\begin{align}
&\frac{N_c g^2}{4\pi^2}
\frac{\Gamma(d/2-1)}{(d-2 )(-z^2)^{d/2-1}}
\int_0^1 \dd u\, \left(u^{3-d}-u\right) \nn &\quad \times
\left(z_{\alpha} G_{z\mu}( \bar uz) - z_{\mu} G_{z\alpha}(\bar uz)\right) \ ,
\label{2a1}
\end{align}
where $G_{z\sigma} \equiv z^\rho G_{\rho\sigma}$ and $\bar u \equiv 1-u$.
The overall
\mbox{$d$-dependent} factor here is finite for $d=4$, but the \mbox{$u$-integral} diverges at the lower limit.
Thus, just like in the case of the link self-energy diagram, the divergence
appears in the integral over a dimensionless parameter $t$ specifying the location
of the endpoint of the gluon line on the gauge link.
The divergence disappears if one uses the UV regularization by taking
$d=4-2\varepsilon_{\rm UV}$, which converts it
into a pole at $\euv=0$.

Since the UV divergence comes from the \mbox{$u \to 0$} integration, we
can isolate it by taking $\bar u =1$ in the gluonic field, which gives
\begin{align}
&\frac{N_c g^2}{4\pi^2}
\frac{\Gamma(d/2-1)}{(d-2 )(-z^2)^{d/2-1}}
\left (\frac1{4-d}-\frac12 \right ) \nn &\quad \times
\left(z_{\alpha} G_{z\mu}(z) - z_{\mu} G_{z\alpha}(z)\right) \ .
\label{UVvert}
\end{align}
The remainder is given by
\begin{align}
&\frac{N_c g^2}{4\pi^2}
\frac{\Gamma(d/2-1)}{(d-2 )(-z^2)^{d/2-1}}
\int_0^1 \dd u\, \left[u^{3-d}-u\right]_{+(0)} \nn &\quad \times
\left(z_{\alpha} G_{z\mu}( \bar uz) - z_{\mu} G_{z\alpha}(\bar uz)\right) \ ,
\label{UVvertreg}
\end{align}
where the plus-prescription at $u=0$ is defined as
\begin{align}
& \int_0^1 \dd u \left[f(u)\right]_{+(0)} g(u) = \int_0^1 \dd u f(u) [g(u) -g(0)] \ .
\label{plus0}
\end{align}

At first sight, the field ${\cal G}_{\mu\alpha}(z)=
z_{\alpha} G_{z\mu}( z) - z_{\mu} G_{z\alpha}(z)$
accompanying the UV pole in Eq. (\ref{UVvert}) does not look like
the field
$G_{\mu\alpha}(z)$ in the original operator.
Thus, one may worry that we are not dealing here with a
multiplicative UV renormalization.
So, let us perform an explicit check for our particular case when
$z= \{0,0,0,z_3 \}$.

To begin with, we see that ${\cal G}_{\mu\alpha}(z)=0$ when both $\mu$ and $\alpha$
are transverse indices $i,j$. This corresponds to a multiplicative renormalization
with the anomalous dimension (AD) equal to zero.

Take now
$\mu=0$. Then ${\cal G}_{0\alpha}(z)= z_{\alpha} G_{z 0}( z)$, so that
${\cal G}_{0 i}(z)= 0$ while ${\cal G}_{0 3}(z)= -z_3^2 G_{3 0}( z) = z_3^2 G_{0 3}( z)$.
Finally, if $\mu=3$, then ${\cal G}_{3\alpha}(z)= - z_{3} G_{z \alpha}( z) = z_3^2 G_{3 \alpha }( z)$,
which gives ${\cal G}_{3 i}(z)= z_3^2 G_{3 i}( z)$ and ${\cal G}_{3 0}(z)= z_3^2 G_{3 0}( z)$
(same result as above).

Thus, for all the cases, ${\cal G}_{\mu\alpha}(z)$ is a multiple of
${G}_{\mu\alpha}(z)$. Namely, when one of the
indices equals 3,
we have a nontrivial anomalous dimension, since
${\cal G}_{3\alpha}(z)=- {\cal G}_{\alpha 3}(z)= z_3^2 G_{3 \alpha}( z)$.
In all other cases, we have a trivial (vanishing) AD, since ${\cal G}_{ij}(z)=0$ and ${\cal G}_{0i}(z)=0$.

As mentioned, the link self-energy diagram has both linear and logarithmic
UV divergences, while the vertex diagrams have just logarithmic
UV divergences.
Adding the logarithmic UV divergence coming from the link self-energy
to the UV divergences of the vertex diagrams, we find, in particular, that
the matrix elements ${ M}_{0 i; i 0}$ and ${ M}_{i j; i j}$ have the logarithmic AD due to the link self-energy diagram only.
Call it $\gamma$.
Comparing overall factors in Eqs. (\ref{selfAD}) and (\ref{UVvert}), we conclude that
${ M}_{3 i; i 3}$ has the logarithmic AD equal to $2 \gamma$ and matrix elements
\mbox{${ M}_{0 i; i 3} \pm
{ M}_{3 i; i 0}$} have the logarithmic AD equal to $\frac32 \gamma$.
In addition, all of these structures acquire at one loop the same factor due to the linear UV singularity.

\subsection{Evolution term}

Our calculations show that the second, UV finite term from the diagram \ref{link}a is given by
\begin{align}
\quad \frac{N_c g^2}{8\pi^2} \frac{\Gamma(d/2-2)}{({d-3})(-z^2)^{d/2-2}} &
\int_0^1 \dd u \left[u^{3-d}-1\right]_{+(0)} \nonumber
\\ &\quad \times
G_{\mu\alpha}( \bar uz) G_{\lambda \beta }( 0) \ .
\label{2aEv}
\end{align}
Note that the gluonic operator in Eq. (\ref{2aEv})
has the same tensor structure as the original operator
$G_{\mu\alpha}( z) G_{\beta \nu}( 0)$ differing from it just by
rescaling $z \to \bar u z$.
There is no mixing with operators of a different type.
The $u$-integral in this case does not diverge for $d=4$, but the overall
\mbox{$\Gamma(d/2-2)$} factor has a pole $1/(d-4)$.

Formally, there is also a pole $1/(d-3)$,
corresponding to a linear UV divergence.
However, the singularity
for $d=3$ is eliminated by the $\left [u^{3-d}-1\right ]$ combination in the integrand.
One may say that the linear divergences present in ``$u^{3-d}$'' and ``$-1$''
parts cancel each other.

In the calculation of Refs.\cite{Zhang:2018diq,Wang:2019tgg} performed using the usual Feynman gauge,
the linear singularities cancel between contributions
of two different diagrams shown in Fig. 1 of Ref. \cite{Wang:2019tgg}.
In our calculation, based on the bF gauge,
the sum of these diagrams is represented by just one vertex diagram,
so that the cancellation
occurs inside the contribution (\ref{2aEv}) of that diagram.

The remaining $1/(d-4)$ pole
corresponds to a collinear divergence
developed
because all the propagators correspond to massless particles.
Taking a nonzero gluon mass $\lambda$, one would get a finite result
containing $K_0(\sqrt{-z^2} \lambda)$
(see, e.g., Ref. \cite{Radyushkin:2017lvu} for a discussion of
the quark vertex diagram in a similar context).

Still, $K_0(\sqrt{-z^2} \lambda)$ is only finite as far as $z^2$ is finite.
The IR cut-off does not eliminate
the logarithmic singularity $\ln (-z^2 \lambda^2)$ that
$K_0(\sqrt{-z^2} \lambda)$ has on the light cone.
In the $z=z_3$ case,
$z_3^2$ works like an ultraviolet cut-off
for this singularity.
This may be contrasted with the UV divergent contributions,
where the UV cut-off is provided by the Polyakov regularization parameter $a$
(or lattice spacing $a_L$) while $z_3^2$ appears on the IR side
of the relevant logarithm $\ln (z_3^2/a^2)$.
